% Follows the supplementary title in X_paired_ablation.
\section{Competitor Diversity and Weighting}
\label{sec:supp_competitor}

We expand the study in \cref{sec:competitor_ablation}. The Core and
Extended collections correspond to the saved \texttt{core} and
\texttt{ext2} menus. These experiments refit saved paired contrasts;
they do not require reconstructing the original intervention triples.
Consequently, the Extended results do not complete the missing
paired-versus-single-sided replication described in
\cref{sec:supp_paired}.

\subsection{Estimators and Sampling Controls}

Let \(\mathcal Q_c\) be the observed competitors of category \(c\),
\(n_{cq}\) the accepted pair count, and \(\bar d_{cq}\) the raw
mean contrast oriented toward \(c\). We define
\begin{equation}
    \widetilde g_c^{(\gamma)}=
    \frac{\sum_{q\in\mathcal Q_c}n_{cq}^{\gamma}\bar d_{cq}}
         {\sum_{q\in\mathcal Q_c}n_{cq}^{\gamma}},
    \qquad
    g_c^{(\gamma)}=
    \frac{\widetilde g_c^{(\gamma)}}{\|\widetilde g_c^{(\gamma)}\|_2}.
    \label{eq:competitor_weighting}
\end{equation}
A0 uses \(\gamma=1\), recovering equal record weights; A1 uses
\(\gamma=0\), giving each observed competitor equal weight; and A2
uses \(\gamma=\tfrac12\). We do not normalize individual contrasts
or pair means. The low-count sensitivity check excludes pairs with fewer
than 10 records from both A0 and A1 before refitting.

For the competitor-count study (B), every eligible category must have
at least eight competitors with at least 16 records each. We sample a
competitor ordering per category and seed and use its first 2, 4, or 8
entries. Each selected competitor contributes 16, 8, or 4 records,
respectively. A shared permutation of each unordered pair's records
couples the fits. The competitor subsets are nested; the complete
32-record sets need not be nested. Records are sampled without replacement
within each category fit, and selections for different categories can
share a record.

The family study (C) requires at least two eligible competitors in the
target's family and two outside it. We compare two same-family
competitors, two cross-family competitors, and one of each, always using
16 records per competitor. Families follow the frozen category-to-family
mapping, including distinct families for the OpenImages categories.
The 21-category Extended cohort contains 15 animal categories; its
composition limits generalization to the full menu.

Area controls preserve 32 records while fixing counts across the three
saved target-area levels to \((8,8,16)\) per category. Each eligible
pair must contain at least \((4,4,8)\) records at those levels. With
\(K\) competitors, each contributes \((8/K,8/K,16/K)\) records.
These quotas control target-area strata, not realized foreground-mask
area or all donor and background properties. Natural-area comparisons
are also evaluated on each controlled cohort to separate cohort changes
from the area intervention.

% Generated by scripts/render_competitor_ablation.py; do not edit numbers.
\begin{table}[t]
    \centering
    \footnotesize
    \setlength{\tabcolsep}{4pt}
    \begin{tabular}{llrr}
        \toprule
        Menu & Weights & AUC $\uparrow$ & AP $\uparrow$ \\
        \midrule
        Core & A0 & 0.9604 & 0.7987 \\
        Core & A1 & 0.9574 & 0.7973 \\
        Core & A2 & 0.9599 & 0.8001 \\
        Extended & A0 & 0.8871 & 0.7042 \\
        Extended & A1 & 0.8885 & 0.7084 \\
        Extended & A2 & 0.8880 & 0.7068 \\
        \bottomrule
    \end{tabular}
    \caption{\textbf{All-record weighting at L18.} Core uses 10 categories and Extended uses 76. A0 weights records equally, A1 weights observed competitors equally, and A2 uses intermediate count weighting ($\gamma=1,0,1/2$ in \cref{eq:competitor_weighting}). Pair means retain their magnitudes; only the final direction is normalized. Metrics use the 0--1 scale.}
    \label{tab:competitor_weighting}
\end{table}

% Generated by scripts/render_competitor_ablation.py; do not edit numbers.
\begin{table*}[t]
    \centering
    \footnotesize
    \setlength{\tabcolsep}{4pt}
    \begin{tabular}{lrrrr}
        \toprule
        Comparison & Categories & $\Delta$ (pp) & 95\% CI (pp) & 97.5\% CI (pp) \\
        \midrule
        A1$-$A0 & 76 & $+0.139$ & $[+0.082,+0.199]$ & $[+0.075,+0.210]$ \\
        $K=8$ minus $K=2$ & 64 & $+1.379$ & $[+1.286,+1.468]$ & $[+1.274,+1.482]$ \\
        \bottomrule
    \end{tabular}
    \caption{\textbf{Primary comparisons on Extended at L18.} The two Bonferroni-adjusted 97.5\% intervals give at least 95\% simultaneous coverage for these comparisons. All differences and endpoints are absolute AUC percentage points. The 2,000 paired bootstrap draws condition on the saved contrasts and fixed fitting seeds.}
    \label{tab:competitor_primary}
\end{table*}


\subsection{Evaluation and Uncertainty}

Core uses the 5,000 COCO val2017 images. Extended evaluates 44 categories
on COCO and 32 additional categories on 25,099 OpenImages train images,
with each category scored only on its assigned source dataset. COCO
instance annotations define presence. For OpenImages, only human-verified
positive or negative image--category labels enter the metrics. We compute
each category's AUC and AP on its eligible images and then average over
the fixed category cohort; the two source datasets are not given equal
weight independently of their category counts.

B and C use fitting seeds 61000--61019. Reported means and sample SDs
refer to per-seed macro metrics, without averaging directions or image
scores. Every layer includes 25 shared random-direction controls.
Image-bootstrap draws resample COCO and OpenImages independently, sharing
each draw across categories, variants, and fitting seeds. We recompute
category AUCs, average over categories and the fixed fits, and form paired
differences within each draw. All reported comparisons retain all 2,000
draws. AP is reported as a point estimate.

The two comparisons specified as primary before this run are Extended
L18 A1 minus A0 and B \(K=8\) minus \(K=2\). We report marginal
95\% intervals and Bonferroni-adjusted 97.5\% intervals, providing at
least 95\% simultaneous coverage for this pair of comparisons
(\cref{tab:competitor_primary}). All other intervals are unadjusted and
secondary. The protocol was defined after inspecting earlier validation
results; it does not constitute evaluation on a previously unseen test set.

These intervals capture image uncertainty conditional on the saved
contrasts and fixed fits. They exclude uncertainty from new donor,
background, gating, or fitting-seed draws. For B at L18, the paired
\(K=8\) minus \(K=2\) seed-difference SD is 0.814 points on Core
and 1.235 points on Extended. These SDs describe variability among fits
and should not be confused with the narrower image-bootstrap intervals.

% Generated by scripts/render_competitor_ablation.py; do not edit numbers.
\begin{table*}[t]
    \centering
    \footnotesize
    \setlength{\tabcolsep}{4pt}
    \begin{tabular}{rrrrr}
        \toprule
        Layer & Core $\Delta$ & Core 95\% CI & Extended $\Delta$ & Extended 95\% CI \\
        \midrule
        1 & $-0.236$ & $[-0.321,-0.150]$ & $+0.132$ & $[+0.078,+0.189]$ \\
        9 & $-0.215$ & $[-0.285,-0.138]$ & $+0.081$ & $[+0.031,+0.135]$ \\
        18 & $-0.298$ & $[-0.376,-0.224]$ & $+0.139$ & $[+0.082,+0.199]$ \\
        27 & $-0.208$ & $[-0.295,-0.119]$ & $+0.036$ & $[-0.024,+0.094]$ \\
        35 & $-0.333$ & $[-0.425,-0.246]$ & $+0.348$ & $[+0.258,+0.437]$ \\
        \bottomrule
    \end{tabular}
    \caption{\textbf{A1 minus A0 across depth.} Differences and intervals use AUC percentage points. Cohorts are the full 10-category Core and 76-category Extended menus. L18 is primary for Extended; the other comparisons are secondary, with unadjusted 95\% intervals.}
    \label{tab:competitor_layers}
\end{table*}

% Generated by scripts/render_competitor_ablation.py; do not edit numbers.
\begin{table*}[t]
    \centering
    \footnotesize
    \setlength{\tabcolsep}{4pt}
    \begin{tabular}{lrrrrrr}
        \toprule
        Menu & Layer & $K=2$ & $K=4$ & $K=8$ & $\Delta_{8-2}$ (pp) & 95\% CI (pp) \\
        \midrule
        Core & 9 & 0.9563 & 0.9629 & 0.9659 & $+0.958$ & $[+0.853,+1.064]$ \\
        Extended & 9 & 0.8438 & 0.8538 & 0.8590 & $+1.513$ & $[+1.426,+1.605]$ \\
        Core & 18 & 0.9610 & 0.9668 & 0.9691 & $+0.811$ & $[+0.718,+0.913]$ \\
        Extended & 18 & 0.8426 & 0.8524 & 0.8564 & $+1.379$ & $[+1.286,+1.468]$ \\
        \bottomrule
    \end{tabular}
    \caption{\textbf{Competitor count at L9 and L18.} Mean macro AUC on the same 7 Core or 64 Extended categories across $K$ and layers, using 20 fits and 32 records per category. AUCs use the 0--1 scale; paired differences use percentage points. Intervals are marginal 95\% image-bootstrap intervals.}
    \label{tab:competitor_budget_layers}
\end{table*}

% Generated by scripts/render_competitor_ablation.py; do not edit numbers.
\begin{table}[t]
    \centering
    \footnotesize
    \setlength{\tabcolsep}{4pt}
    \begin{tabular}{lrrr}
        \toprule
        Area setting & $K$ & AUC $\uparrow$ & AP $\uparrow$ \\
        \midrule
        Natural & 2 & $0.8366\pm0.0131$ & 0.6076 \\
        Natural & 4 & $0.8470\pm0.0108$ & 0.6321 \\
        Natural & 8 & $0.8516\pm0.0078$ & 0.6423 \\
        Controlled & 2 & $0.8470\pm0.0111$ & 0.6368 \\
        Controlled & 4 & $0.8564\pm0.0095$ & 0.6549 \\
        Controlled & 8 & $0.8638\pm0.0081$ & 0.6710 \\
        \bottomrule
    \end{tabular}
    \caption{\textbf{Area control on Extended at L18.} Both settings use the same 54 eligible categories, 32 records per category, and 20 fits. Controlled fits use per-category area-level counts $(8,8,16)$. AUC is mean $\pm$ seed SD; AP is a mean, both on the 0--1 scale. Core has no eligible category for this control.}
    \label{tab:competitor_area}
\end{table}

% Generated by scripts/render_competitor_ablation.py; do not edit numbers.
\begin{table*}[t]
    \centering
    \footnotesize
    \setlength{\tabcolsep}{4pt}
    \begin{tabular}{lllrr}
        \toprule
        Menu & Area setting & Competitors & AUC $\uparrow$ & AP $\uparrow$ \\
        \midrule
        Core & Natural & Same & $0.9504\pm0.0348$ & 0.8047 \\
        Core & Natural & Mixed & $0.9656\pm0.0146$ & 0.8255 \\
        Core & Natural & Cross & $0.9718\pm0.0029$ & 0.8469 \\
        Core & Controlled & Same & $0.9647\pm0.0153$ & 0.8386 \\
        Core & Controlled & Mixed & $0.9674\pm0.0136$ & 0.8043 \\
        Core & Controlled & Cross & $0.9664\pm0.0074$ & 0.7660 \\
        Extended & Natural & Same & $0.7983\pm0.0266$ & 0.5709 \\
        Extended & Natural & Mixed & $0.8595\pm0.0157$ & 0.6576 \\
        Extended & Natural & Cross & $0.8680\pm0.0139$ & 0.6733 \\
        Extended & Controlled & Same & $0.8165\pm0.0147$ & 0.6017 \\
        Extended & Controlled & Mixed & $0.8664\pm0.0191$ & 0.6947 \\
        Extended & Controlled & Cross & $0.8765\pm0.0139$ & 0.7035 \\
        \bottomrule
    \end{tabular}
    \caption{\textbf{Competitor-family selection at L18.} All fits use two competitors and 32 records per category. Core uses three categories; Extended uses 21, with the same cohorts under area control. AUC is mean $\pm$ seed SD over 20 fits; AP is the mean macro AP (0--1 scale). These comparisons are secondary.}
    \label{tab:competitor_family}
\end{table*}

% Generated by scripts/render_competitor_ablation.py; do not edit numbers.
\begin{table}[t]
    \centering
    \footnotesize
    \setlength{\tabcolsep}{4pt}
    \begin{tabular}{llrr}
        \toprule
        Menu & Area & $\Delta$ (pp) & 95\% CI (pp) \\
        \midrule
        Core & Natural & $+2.136$ & $[+1.417,+2.879]$ \\
        Core & Controlled & $+0.170$ & $[-0.467,+0.787]$ \\
        Extended & Natural & $+6.968$ & $[+6.357,+7.587]$ \\
        Extended & Controlled & $+5.997$ & $[+5.388,+6.621]$ \\
        \bottomrule
    \end{tabular}
    \caption{\textbf{Cross-family minus same-family AUC at L18.} Percentage points with unadjusted 95\% paired image-bootstrap intervals. Intervals condition on the saved contrasts and fixed 20 fits.}
    \label{tab:competitor_family_deltas}
\end{table}

% Generated by scripts/render_competitor_ablation.py; do not edit numbers.
\begin{table*}[t]
    \centering
    \footnotesize
    \setlength{\tabcolsep}{4pt}
    \begin{tabular}{llrp{0.69\textwidth}}
        \toprule
        Menu & Study & $n$ & Eligible categories \\
        \midrule
        Core & B & 7 & car, bus, dog, cat, bird, person, bottle \\
        Core & B-area & 0 & No eligible category \\
        Core & C / C-area & 3 & dog, cat, bird \\
        Extended & B & 64 & car, bus, cat, bird, chair, couch, tv, bicycle, motorcycle, train, traffic light, fire hydrant, stop sign, horse, sheep, cow, elephant, bear, zebra, giraffe, umbrella, handbag, frisbee, skateboard, wine glass, cup, bowl, sandwich, pizza, donut, cake, toilet, laptop, book, clock, vase, Ambulance, Balloon, Camel, Camera, Christmas tree, Cucumber, Dolphin, Goldfish, Guitar, Ipod, Kettle, Lion, Microwave oven, Plastic bag, Pumpkin, Rabbit, Rhinoceros, Shark, Sombrero, Squirrel, Teapot, Tennis ball, Watermelon, Wok, Zucchini, barrel, flag, sculpture \\
        Extended & B-area & 54 & bus, cat, bird, chair, couch, motorcycle, train, fire hydrant, stop sign, horse, sheep, cow, elephant, bear, zebra, giraffe, umbrella, handbag, frisbee, skateboard, wine glass, bowl, sandwich, pizza, donut, toilet, laptop, book, clock, vase, Ambulance, Balloon, Camel, Camera, Christmas tree, Cucumber, Dolphin, Goldfish, Guitar, Ipod, Kettle, Microwave oven, Plastic bag, Pumpkin, Rhinoceros, Shark, Sombrero, Teapot, Tennis ball, Watermelon, Zucchini, barrel, flag, sculpture \\
        Extended & C / C-area & 21 & dog, bird, motorcycle, horse, sheep, cow, elephant, bear, zebra, giraffe, cake, Balloon, Cucumber, Dolphin, Goldfish, Monkey, Rabbit, Rhinoceros, Shark, Tennis ball, Zucchini \\
        \bottomrule
    \end{tabular}
    \caption{\textbf{Exact sampling cohorts.} B requires eight supported competitors; C requires two same-family and two cross-family competitors. Area controls additionally require per-pair support at all three target-area levels. Each cohort is fixed across compared variants and fitting seeds. Names retain the saved menu capitalization.}
    \label{tab:competitor_cohorts}
\end{table*}


\subsection{Secondary Results and Scope}

Equal competitor weights reduce Core AUC at all five layers. Extended
point estimates improve at all five layers, although the L27 interval
includes zero (\cref{tab:competitor_layers}). Removing pairs with fewer
than 10 records attenuates the L18 differences to \(-0.054\) points
on Core and \(+0.064\) on Extended. This sensitivity does not identify
the noise level of those records or justify choosing a cutoff after
observing the result.

The competitor-count gain appears at L9 and L18. Under area control,
Extended \(K=8\) exceeds \(K=2\) by 1.681 points at L18
(95\% CI \([1.580,1.773]\); \cref{tab:competitor_area}). No Core
category satisfies the B-area support constraints. Full-record A0 remains
above the 32-record \(K=8\) fit on the same B cohorts: 0.9735 versus
0.9691 on Core, and 0.8860 versus 0.8564 on Extended. These references
use unequal budgets and do not contradict the fixed-budget result.

Area control preserves the cross-family advantage on Extended, with a
5.997-point L18 difference (95\% CI \([5.388,6.621]\)). On Core,
the interval crosses zero and cross-family AP is lower than same-family
AP (0.7660 versus 0.8386; \cref{tab:competitor_family}). We therefore do
not claim a robust Core advantage or a general semantic-distance effect.

The selected class diagnostics also resist a universal balancing claim.
At Extended L18, A1 changes chair AUC from 0.6100 to 0.6252, couch from
0.6344 to 0.6347, and handbag from 0.4708 to 0.4613. The corresponding
unadjusted class-level intervals are exploratory. Restricting A1 minus A0
to B's 64-category cohort changes its mean effect to \(-0.050\) points,
reinforcing the need to retain exact cohort definitions.

\subsection{Reconstruction and Reproducibility}

We use model revision \texttt{66285546d2b8}, prompt
\texttt{Describe.}, BF16 inference, FP16 feature caches, and FP32 cosine
scoring. Layers 1, 9, 18, 27, and 35 are extracted in one forward pass
per image on two H100 80\,GB GPUs. Weighting is evaluated at all five
layers; B, C, and their area controls are evaluated at L9 and L18.
The run records 148 passing regression tests, exact agreement of
extraction hooks with model reference outputs, and exact recovery of
historical A0 directions at every layer in both collections. The maximum
per-category AUC difference between GPU scoring and the original CPU
readout is \(3.59\times10^{-6}\), below the \(10^{-5}\) tolerance.

All variants share a reconstructed evaluation cache. The historical
OpenImages run contained 25,085 images; the identities of the 14 omitted
images and the historical model revision are unavailable. Changes between
historical and reconstructed scores are therefore not ablation effects.
Core A0 also differs slightly from the replayed P0 in
\cref{sec:paired_ablation}, which recomputes contrasts from recovered
triples. Core and Extended differ in their menus, gates, and saved
contrasts, so their macro scores do not form a controlled menu-size test.

The saved artifacts include exact record selections, category-family
snapshots, eligibility audits, input hashes, per-image scores, and
bootstrap samples. A larger-budget study with 64 records per category
and broader family coverage remain future experiments; no results from
those proposed studies are included here.
